% Appendix A: Source Code and Setup

\section{Source Code Repository}

The complete source code for the Faculty Feedback Management System is available on GitHub:

\begin{center}
\url{https://github.com/Huzaif2309/faculty-feedback-system}
\end{center}

The repository contains the full implementation including the frontend application, backend API, and database schema.

\section{Frontend Setup}

\subsection{Prerequisites}

\begin{itemize}
    \item Node.js (v18 or higher)
    \item Supabase Account (\url{https://supabase.com})
    \item Google Cloud Console Project for OAuth and AI features
    \item Resend Account for email functionality (\url{https://resend.com})
\end{itemize}

\subsection{Installation}

\begin{lstlisting}[language=bash, caption={Clone and Install Frontend}]
git clone https://github.com/Huzaif2309/faculty-feedback-system
cd faculty-feedback-system/frontend
npm install
\end{lstlisting}

\subsection{Environment Configuration}

\begin{lstlisting}[language=bash, caption={Configure Environment Variables}]
cp env.example .env
\end{lstlisting}

Required environment variables in \texttt{.env}:
\begin{itemize}
    \item \texttt{NEXT\_PUBLIC\_SUPABASE\_URL} - Supabase project URL
    \item \texttt{NEXT\_PUBLIC\_SUPABASE\_ANON\_KEY} - Supabase anonymous key
    \item \texttt{RESEND\_API\_KEY} - Resend API key for emails
    \item \texttt{GOOGLE\_GENERATIVE\_AI\_API\_KEY} - Google AI API key
    \item \texttt{NEXT\_PUBLIC\_VAPID\_PUBLIC\_KEY} - VAPID public key
    \item \texttt{VAPID\_PRIVATE\_KEY} - VAPID private key
\end{itemize}

Generate VAPID keys for push notifications:

\begin{lstlisting}[language=bash, caption={Generate VAPID Keys}]
npx web-push generate-vapid-keys
\end{lstlisting}

\subsection{Database Setup}

Initialize the Supabase database by running the SQL schema in the Supabase Dashboard SQL Editor:

\begin{lstlisting}[language=bash, caption={Database Schema Location}]
frontend/lib/supabase/migrations/schema.sql
\end{lstlisting}

\subsection{Run the Application}

\begin{lstlisting}[language=bash, caption={Start Frontend Development Server}]
npm run dev
# Application available at http://localhost:3000
\end{lstlisting}

\section{Backend Setup}

The backend is a Flask API that provides additional REST endpoints for data management.

\subsection{Prerequisites}

\begin{itemize}
    \item Python 3.10 or higher
    \item pip package manager
\end{itemize}

\subsection{Installation}

\begin{lstlisting}[language=bash, caption={Install Backend Dependencies}]
cd faculty-feedback-system/faculty-feedback-backend
pip install -r requirements.txt
\end{lstlisting}

\subsection{Environment Configuration}

Create a \texttt{.env} file with the following variables:

\begin{itemize}
    \item \texttt{SUPABASE\_URL} - Supabase project URL
    \item \texttt{SUPABASE\_KEY} - Supabase service role key
\end{itemize}

\subsection{Run the Backend}

\begin{lstlisting}[language=bash, caption={Start Backend Server}]
python app.py
# API available at http://localhost:5000
\end{lstlisting}

\section{Docker Deployment}

Both frontend and backend can be deployed using Docker:

\begin{lstlisting}[language=bash, caption={Docker Deployment}]
docker-compose up --build
\end{lstlisting}


